\chapter{Einleitung} \label{kap:1}

Digitale Anwendungen sind für Verwaltung, Wirtschaft und Gesellschaft zu einer zentralen Infrastruktur geworden. Hinsichtlich dieser Entwicklung gewinnt Web-Barrierefreiheit zunehmend an Bedeutung, damit alle Menschen - unabhängig von individuellen Einschränkungen - gleichberechtigten Zugang zu digitalen Informationen und Dienstleistungen erhalten. In der Praxis zeigt sich jedoch eine Diskrepanz zwischen normativen Anforderungen (z. B. \ac{WCAG}-basierte Kriterien)\footcite[Vgl.][]{WCAG} und der tatsächlichen Umsetzung in Softwareprojekten: Barrierefreiheit wird häufig spät berücksichtigt, Prüfungen erfolgen punktuell und Korrekturen sind zeitaufwendig. Besonders in gewachsenen Systemlandschaften mit Legacy-Anteilen und knappen Entwicklungsressourcen entstehen so immer wieder Qualitätsdefizite.

Parallel dazu haben sich \ac{LLM}s als leistungsfähige Werkzeuge zur Unterstützung von Softwareentwicklung etabliert. Sie helfen, Aufgaben wie Code-Analyse, Refactoring oder die Erstellung technischer Empfehlungen zu beschleunigen. Für die Barrierefreiheit ist dies besonders relevant: Hier müssen nicht nur Fehler erkannt, sondern auch korrekt, standardkonform und kontextsensitiv behoben werden. Gleichzeitig bestehen in vielen Organisationen, insbesondere in sensiblen Umgebungen, strikte Vorgaben an Datenschutz und Informationssicherheit, wodurch Cloud-basierte Assistenzsysteme häufig nicht eingesetzt werden dürfen. Daraus ergibt sich die zentrale Herausforderung, \ac{LLM}-Unterstützung so zu gestalten, dass sie lokal, nachvollziehbar und qualitätsgesichert in bestehende Entwicklungsprozesse integriert werden kann.

Diese Arbeit untersucht daher, wie ein \ac{LLM}- bzw. agentenbasierter Ansatz zur automatisierten Erkennung und Behebung von Barrierefreiheitsproblemen in Webanwendungen konzipiert und prototypisch umgesetzt werden kann.\footcite[Vgl.][S. 297]{bassi_assessment_2025} Im Fokus steht ein Workflow, der automatisierte Tests (z. B. via Browser-Automation und Accessibility-Linter) mit einer \ac{LLM}-gestützten Analyse und Änderungsvorschlägen verbindet und die Änderungen anschließend erneut testet. Ziel ist es, den manuellen Aufwand zu reduzieren, die Konsistenz der Ergebnisse zu erhöhen und Barrierefreiheit als kontinuierlichen Bestandteil der Softwarequalität zu verankern.

\section{Motivation und Kontext}

Rund 1,3 Milliarden Menschen (16\% der Weltbevölkerung) leben mit einer Form von Behinderung.\footcite[Vgl.][]{WHO} Für diese Personengruppe ist der Zugang zu digitalen Angeboten keine bloße Komfortfrage, sondern eine Voraussetzung für gesellschaftliche Teilhabe in nahezu allen Lebensbereichen: von der Kommunikation mit Behörden über den Zugang zu Bildung und Gesundheitsinformationen bis hin zur Teilnahme am wirtschaftlichen Leben. Das Recht auf barrierefreien Zugang zu Informations- und Kommunikationstechnologien ist völkerrechtlich in der UN-Behindertenrechtskonvention verankert und hat sich in den vergangenen Jahren zunehmend auch in verbindlichem nationalen und europäischen Recht niedergeschlagen.\footcite[Vgl.][]{OHCHR}

Auf europäischer Ebene verpflichtet der \ac{EAA} die Mitgliedstaaten, Barrierefreiheitsanforderungen für digitale Produkte und Dienstleistungen durchzusetzen.\footcite[Vgl.][]{EAA} In Deutschland wurde diese Richtlinie durch das \ac{BFSG} in nationales Recht überführt, das seit dem 28. Juni 2025 auch für weite Teile der Privatwirtschaft gilt.\footcite[Vgl.][]{BFSG} Ergänzt wird dieser Rahmen durch die \ac{BITV}, die für Behörden und öffentliche Stellen bereits seit Jahren verbindliche Vorgaben formuliert.\footcite[Vgl.][S. 281]{BITV, heinemann_barrierefreiheit_2024}

Trotz dieser eindeutigen Rechtslage ist die praktische Umsetzung barrierefreier Webanwendungen nach wie vor unzureichend. Die jährlich durchgeführte WebAIM-Million-Studie, die die Startseiten der meistbesuchten Webseiten weltweit analysiert, stellte 2025 fest, dass 94,8\% der untersuchten Seiten nachweisbare Verstöße gegen die \ac{WCAG} aufweisen.\footcite[Vgl.][]{WebAIM} Damit bleibt Barrierefreiheit in der Entwicklungspraxis trotz gesetzlicher Verpflichtung systematisch unterrepräsentiert. Im Vergleich mit den letzten Jahren ist seit 2019 (97,8\%) lediglich eine Verbesserung um 3 Prozentpunkte zu verzeichnen.\footcite[Vgl.][]{WebAIM19}

Die \ac{BITBW} nimmt als zentrale IT-Dienstleisterin der Landesverwaltung Baden-Württemberg in diesem Kontext eine besondere Stellung ein.\footcite[Vgl.][]{BITBW_2021} Als Entwicklerin und Betreiberin öffentlicher Webanwendungen für Ministerien, Polizeibehörden und Kommunen ist sie gesetzlich zur Einhaltung der \ac{BITV} und des \ac{BFSG} verpflichtet.\footcite[Vgl.][]{Landesrecht_BW} Drei Faktoren machen die \ac{BITBW} zu einem besonders geeigneten Praxiskontext für die vorliegende Arbeit: Erstens unterliegen ihre Anwendungen strengeren Barrierefreiheitsanforderungen als privatwirtschaftliche Produkte. Zweitens schließen die Datenschutzvorgaben der öffentlichen Verwaltung die Nutzung cloudbasierter \ac{LLM}-\ac{API}s weitgehend aus. Drittens erschweren gewachsene Systemlandschaften mit Legacy-Anteilen und begrenzten Entwicklungsressourcen die manuelle Barrierefreiheitsprüfung im erforderlichen Umfang. Sie bildet damit den praktischen Rahmen der vorliegenden Arbeit.

\section{Problemstellung}

Die mangelhafte Umsetzung von Barrierefreiheit in Softwareprojekten lässt sich nicht allein auf fehlende Motivation oder Ressourcen zurückführen. Ein wesentlicher Treiber sind die Schwächen der verfügbaren Werkzeuge und Tools: Automatisierte Barrierefreiheitsanalyse-Tools wie axe\footcite[Vgl.][]{axe}, WAVE\footcite[Vgl.][]{WAVE} oder Lighthouse\footcite[Vgl.][]{Lighthouse} erkennen zwar syntaktische Verstöße, stoßen jedoch bei semantischen Problemen an ihre Grenzen. Ob ein vorhandener Alternativtext den Bildinhalt tatsächlich angemessen beschreibt, oder ob eine Schaltflächenbeschriftung für Screenreader-Nutzerinnen und -Nutzer verständlich ist, entzieht sich der regelbasierten Prüflogik dieser Werkzeuge.

Empirisch belegt wird diese Schwäche durch das Mutation-Testing-Framework \textfett{Ma11y}: In einer systematischen Evaluation mit 25 gezielt eingebrachten Barrierefreiheits-Defekten erkannte kein einzelnes der sechs untersuchten Tools mehr als 50\% der injizierten Fehler. Somit blieben im Durchschnitt nahezu die Hälfte aller Defekte unentdeckt.\footcite[Vgl.][S. 100]{tafreshipour_ma11y_2024} Besonders aufschlussreich ist dabei die Differenzierung nach Verstoßtypen: Während syntaktische Defekte zu 54\% erkannt wurden, lag die Erkennungsrate bei semantischen Verstößen lediglich bei 26\%.\footcite[Vgl.][S. 108]{tafreshipour_ma11y_2024} Sechs der 25 Defekttypen - darunter die Manipulation von Seitentiteln und Alternativtexten - wurden von keinem der untersuchten Tools erkannt.

Eine zweite strukturelle Schwäche betrifft die Charakteristik moderner Webanwendungen. Rich-Client-Anwendungen auf Basis von Frameworks wie React erzeugen ihren Inhalt größtenteils clientseitig und verändern den \ac{DOM} dynamisch in Abhängigkeit von Nutzerinteraktionen. Die zunehmende Komplexität solcher Webanwendungen ist im Anhang veranschaulicht (vgl. Abbildung \ref{fig:complexity}). Barrieren entstehen häufig erst in bestimmten Zuständen: wenn ein Modalfenster geöffnet wird, eine Fehlermeldung erscheint oder ein Dropdown-Menü aktiviert ist. Rein statische Analysen - oder Analysen, die ausschließlich den initialen Seitenladezustand prüfen - erfassen diese zustandsabhängigen Barrieren nicht.\footcite[Vgl.][S. 303]{bassi_assessment_2025}

Jüngere Forschungsarbeiten zeigen, dass \ac{LLM}s geeignet sind, die Grenzen regelbasierter Tools zu überwinden. So erreicht das System \textfett{AccessGuru}durch den Einsatz von GPT-4o eine Reduktion des Violation Scores um bis zu 84\%\footcite[Vgl.][]{fathallah_accessguru_2025}, und \textfett{GenA11y} erzielt bei der Erkennung von Barrierefreiheitsproblemen eine Precision von 94,5\%\footcite[Vgl.][FSE101:1]{he_enhancing_2025}. Diese Ansätze beschränken sich jedoch auf statisches \ac{HTML}. Keiner der bekannten Ansätze verknüpft Tool-Reports, Playwright-\allowbreak Interaktionsdaten und den Quellcode der geprüften Komponenten zu einem integrierten Analysesystem.\footcite[Vgl.][]{fathallah_accessguru_2025} Der Schritt von der Fehlererkennung zu einer entwicklernahen, direkt umsetzbaren Handlungsempfehlung - idealerweise mit Bezug zur konkreten Quellcodestelle - wird in der Literatur zwar als offenes Problem benannt,\footcite[Vgl.][FSE101:15-17]{he_enhancing_2025} bleibt aber bislang ungelöst.

Daraus ergibt sich eine klar umrissene Lücke: Es fehlt ein Assistenzsystem, das (a) die dynamischen Zustände einer React-Anwendung durch automatisierte Browser-Interaktionen erfasst, (b) die Ergebnisse automatisierter Tool-Reports mit dem relevanten Quellcode verknüpft und (c) auf dieser Grundlage priorisierte, entwicklernahe Handlungsempfehlungen erzeugt - lokal ausführbar und ohne den Transfer sensibler Anwendungsdaten an externe Dienste.

\section{Zielsetzung und Forschungsfragen} \label{kap:1.3}

Ziel dieser Bachelorarbeit ist die Konzeption und prototypische Umsetzung einer lokal ausführbaren, kontextsensitiven Assistenzlösung zur Barrierefreiheitsanalyse von Rich-Client-Web\-an\-wen\-dun\-gen. Als Untersuchungsobjekt dient eine React-basierte Webanwendung aus dem Umfeld der \ac{BITBW}. Der Proof of Concept soll automatisierte Barrierefreiheits-Reports, Interaktionsdaten aus Playwright-\allowbreak Tests sowie ausgewählte Informationen aus der Codebasis in einer Datenpipeline zusammenführen. Auf dieser Grundlage werden priorisierte, entwicklernahe Handlungsempfehlungen abgeleitet und strukturiert aufbereitet.

Die Arbeit liefert drei konkrete Beiträge: Erstens einen strukturierten Überblick darüber, wie unterschiedliche Datenquellen der Barrierefreiheitsanalyse systematisch kombiniert und interpretiert werden können. Zweitens einen Proof of Concept, der das Vorgehen exemplarisch demonstriert. Drittens eine qualitative Evaluation anhand definierter Kriterien, die aus den Anforderungen und der einschlägigen Fachliteratur abgeleitet werden.

Die übergeordnete Forschungsfrage dieser Arbeit lautet:

\textfett{Inwieweit kann ein kontextsensitives KI-Assistenzsystem, das Tool-Reports, Play\-wright-Inter\-aktions\-daten und Quellcode kombiniert, Barrierefreiheitsverstöße in Re\-act-ba\-sier\-ten Webanwendungen automatisiert erkennen und entwicklernahe Handlungsempfehlungen generieren?}

Diese übergeordnete Frage wird durch zwei konkretisierende \ac{TF} operationalisiert:
\begin{itemize}
    \item \ac{TF}1: Welche Typen von Barrierefreiheitsverstößen - syntaktisch, semantisch und layout-bezogen - lassen sich durch den gewählten Ansatz zuverlässig erkennen, und wo liegen die methodischen Grenzen?
    \item \ac{TF}2: Inwiefern verbessert die Einbeziehung von Quellcode und Playwright-\allowbreak Interaktionsdaten die Qualität der generierten Handlungsempfehlungen gegenüber rein \ac{DOM}-basierten Ansätzen?
\end{itemize}

\ac{TF}1 adressiert dabei das Erkennungsproblem und erlaubt eine differenzierte Aussage über die Leistungsfähigkeit und Grenzen des Systems entlang der Verstoß-Taxonomie von Fathallah et al.\footcite[Vgl.][S. 114-115]{fathallah_accessguru_2025, Šmite_Wohlin_Galviņa_Prikladnicki_2012} \ac{TF}2 greift den in der Literatur identifizierten Mehrwert der Quellcode-Kontextualisierung auf und überprüft ihn im Rahmen einer qualitativen Evaluation am gewählten Untersuchungsobjekt.

\section{Aufbau der Arbeit}

Die vorliegende Arbeit gliedert sich in acht Kapitel. Im Anschluss an die Einleitung werden in Kapitel \ref{kap:2} die theoretischen Grundlagen erarbeitet: Barrierefreiheit im Web, die \ac{WCAG}-Standards und der rechtliche Rahmen, die Charakteristika von Rich-Client-Web\-an\-wen\-dun\-gen sowie Grundlagen zu Large Language Models und Prompt-Engineering-Strategien.

Kapitel \ref{kap:3} gibt einen systematischen Überblick über den Stand der Forschung. Die relevanten Arbeiten werden thematisch gegliedert - nach \ac{LLM}-basierter Erkennung, Code-Korrektur, hybriden Pipelines und Evaluation - und in einer Konzeptmatrix nach Webster und Watson zusammengefasst.\footcite[Vgl.][]{WebsterWatson2002} Das Kapitel schließt mit einer expliziten Einordnung des Beitrags der vorliegenden Arbeit.

Kapitel \ref{kap:methodik} legt das methodische Vorgehen dar. Es beschreibt die Wahl des Design-Science-Research-Paradigmas als wissenschaftliche Methode, definiert den Untersuchungsgegenstand und entwirft das Evaluationsdesign.\footcite[Vgl.][]{hevner2004dsr} Kapitel \ref{kap:4} analysiert den Problemkontext und leitet daraus funktionale und nicht-funktionale Anforderungen ab, einschließlich der Begründung der Modellauswahl.

Auf Basis dieser Anforderungen wird in Kapitel \ref{kap:5} das Assistenzsystem konzipiert: Systemarchitektur, Datenpipeline, Prompt-Strategie und Ausgabeformat werden begründet entworfen. Kapitel \ref{kap:6} beschreibt die prototypische Implementierung des Proof of Concept und illustriert die Systemfunktionalität anhand von Beispielabläufen für syntaktische, semantische und dynamisch erzeugte Barrierefreiheitsprobleme.

Kapitel \ref{kap:7} präsentiert und diskutiert die Evaluationsergebnisse, beantwortet die Forschungsfragen und reflektiert die Limitationen des Ansatzes. Das abschließende Kapitel \ref{kap:8} fasst die Kernaussagen zusammen, ordnet den Beitrag der Arbeit in die Forschungslandschaft ein und formuliert Empfehlungen für weiterführende Untersuchungen.